\hypertarget{appendix-02b9d9a81dca}{}
\section*{Barycentric coordinates and interpolation}
\section{Barycentric coordinates and interpolation}\label{barycentric-coordinates-and-interpolation}

\textbf{Curriculum depth:} Developed · \textbf{First seen:} Phase6 · \textbf{Latest development:} University courses

\subsection{Definition and mechanism}\label{definition-and-mechanism}

Barycentric weights reconstruct a triangle point as an affine combination. The weights sum to one; nonnegative weights describe points inside or on a nondegenerate triangle. Dot-system solves on off-plane inputs describe their projection, not proven surface membership.

The panel sensor \((1,1,0)\) has weights \((5/12,1/4,1/3)\) and corner temperatures \((20,40,50)\) interpolate to \(35\).

\subsection{Prerequisites}\label{prerequisites}

\begin{itemize}
\tightlist
\item
  \hyperref[appendix-3b863809f82a]{Gram matrices and coordinate metrics} --- Pairwise basis measurements define the systems and metric used in reconstruction.
\item
  \hyperref[appendix-0659864c0aa3]{Vectors coordinates and affine points} --- Point, displacement and coordinate roles determine which geometric operations are meaningful.
\end{itemize}

\subsection{Used by}\label{used-by}

\begin{itemize}
\item
  \hyperref[appendix-0026ec42c9d5]{Convex combinations and convex hulls} --- Builds on this mechanism in the connected geometry story.
\item
  \href{}{Triangle rasterization and coverage tests (04-Visual-Computing-II, appendix-ed3f9bc298fe)}
\item
  \href{}{Ray-triangle intersection (04-Visual-Computing-II, appendix-47ac1bd6d144)}
\item
  \href{}{Stereo correspondence and epipolar geometry (03-Visual-Computing-I, appendix-c4e98b6e4001)}
\item
  \href{}{Homogeneous clipping and perspective interpolation (03-Visual-Computing-I, appendix-40426e4b2103)}
\item
  \hyperref[appendix-0c2a5f69851e]{Interpolation versus approximation} \#\# Recall and next depth
\end{itemize}

Derive signed-area ratios; distinguish affine attribute interpolation from perspective-correct screen interpolation. Full production proofs, target-specific integration and learner mastery require further evidence.

\subsection{Source and navigation}\label{source-and-navigation}

\href{Continuous\%20Notes/Computing\%20Foundations\%20-\%20From\%20Bits\%20to\%20C++\%20and\%20Systems\%20APIs\#The\%20sensor\%20needs\%20coordinates\%20inside\%20the\%20triangle}{Teaching narrative} · \href{}{Coverage and source gap (01-Foundations-I, appendix-fd090d982a15 § Phase 6 - Parts 501-600)} · \href{}{Evidence and limits (01-Foundations-I, appendix-a103e7846d21 § Phase 6 - Source analysis and verification)} · \hyperref[appendix-4ad47bdd269a]{Dependency Map}

\subsection{University course development}\label{university-course-development}

Extend the existing interpolation note with the geometric meanings that make triangle rasterization and camera calibration easier to reason about. For triangle vertices \(p_0,p_1,p_2\), a point inside can be written \(p=\lambda_0p_0+\lambda_1p_1+\lambda_2p_2\), where the weights sum to one and are nonnegative. Each weight is the area of the opposite subtriangle divided by the full triangle area; equivalently it is a signed edge-distance normalized by altitude. The weights therefore provide both a point-in-triangle test and attribute interpolation. They also reveal a limitation: screen-space weights do not linearly interpolate a 3D attribute after perspective projection. Interpolate reciprocal depth and attribute divided by depth, then divide the interpolated values to recover perspective-correct texture coordinates or colors.

\href{Continuous\%20Notes/University\%20-\%20Worlds,\%20Images\%20and\%20Inference\%20-\%20Graphics,\%20Vision\%20and\%20Photography\#Rasterization,\%20coverage,\%20and\%20depth}{Rasterization, coverage, and depth} · \href{Continuous\%20Notes/University\%20-\%20Worlds,\%20Images\%20and\%20Inference\%20-\%20Graphics,\%20Vision\%20and\%20Photography\#Textures\%20and\%20surface\%20detail}{Textures and surface detail} · \href{Continuous\%20Notes/University\%20-\%20Worlds,\%20Images\%20and\%20Inference\%20-\%20Graphics,\%20Vision\%20and\%20Photography\#Pose\%20estimation,\%20transforms,\%20and\%20robust\%20fitting}{Pose estimation, transforms, and robust fitting}
